% !TeX program = lualatex
% !TeX encoding = utf8
% !TeX spellcheck = uk_UA
% !TeX root = ../main.tex

%% --------------------------------------------------------
%% Chapter file: Kinetics.tex
%% --------------------------------------------------------

\graphicspath{{\currfilebase/Pictures}}



\chapter{Кінетичний опис плазми}\label{sec:plasma-kinetic}\hypertarget{plasma-kinetic}{}
%% --------------------------------------------------------

Рідинний опис виникає після усереднення руху великої кількості частинок. Але іноді самого
середнього значення швидкості недостатньо. Дві плазми можуть мати однакові густину $n$ та
середню швидкість $\vect v$, але різні розподіли частинок за швидкостями, а~отже
по-різному реагувати на хвилі, зіткнення та нестійкості.

Кінетичний опис зберігає інформацію про те, як частинки розподілені не лише в
координатному просторі, а~й~у~просторі швидкостей. Це особливо важливо тоді, коли характер
процесу визначається невеликою частиною частинок із певними швидкостями: наприклад, при
резонансній взаємодії з~хвилею. Рідинні величини в~такому разі вже можуть не містити
достатньої інформації.



% --------------------------------------------------------
\section{Функція розподілу та фазовий простір}
% --------------------------------------------------------



Для кожного сорту частинок $s$ вводять \emphx{функцію розподілу} $f_s(\vect r,\vect v,t)$.
Вона визначена у~шестимірному фазовому просторі $(\vect r,\vect v)$ і~часі. Величина

\begin{equation*}
f_s(\vect r,\vect v,t)\,d^3r\,d^3v
\end{equation*}

є середнім числом частинок виду $s$, які в~момент часу $t$ перебувають у~елементі фазового
простору $d^3r\,d^3v$.

На відміну від концентрації, функція $f_s$ містить інформацію про розподіл частинок за
швидкостями. Концентрація є її нульовим моментом,

\begin{equation}\label{eq:kinetic-density}
 n_s(\vect r,t)=\int f_s(\vect r,\vect v,t)\,d^3v,
\end{equation}

а~середня швидкість визначається першим моментом,

\begin{equation}\label{eq:kinetic-velocity}
 n_s\vect v_s=\int \vect v f_s\,d^3v.
\end{equation}

Знання тільки $n_s$ та $\vect v_s$ не відновлює функцію розподілу. У~цьому полягає одна
з~головних відмінностей кінетичного опису від рідинного: кінетична теорія зберігає значно
більше інформації про мікроскопічний стан плазми.



% --------------------------------------------------------
\section{Моменти функції розподілу}
% --------------------------------------------------------



Середня швидкість характеризує рух центра розподілу. Для опису хаотичного руху введемо
відносну швидкість

\begin{equation*}
\vect c=\vect v-\vect v_s.
\end{equation*}

Тоді тензор тиску визначається другим центральним моментом:

\begin{equation}\label{eq:pressure-tensor}
 P_{s,ij}
 =
 m_s\int
 (v_i-v_{s,i})(v_j-v_{s,j})f_s\,d^3v.
\end{equation}

У~загальному випадку $\mathsf P_s$ є тензором, а~не скаляром. Якщо розподіл ізотропний,
усі напрямки еквівалентні й

\begin{equation}\label{eq:isotropic-pressure}
 P_{s,ij}=p_s\delta_{ij}.
\end{equation}

Скалярний тиск тоді пов'язаний із другим центральним моментом:

\begin{equation}\label{eq:pressure-energy}
 p_s=\frac{m_s}{3}\int c^2 f_s\,d^3v.
\end{equation}

Звідси видно фізичний зміст температури: вона характеризує середню енергію хаотичного руху
частинок. У~неізотропній плазмі природніше говорити окремо про тиск поперек і~вздовж
магнітного поля, $p_\perp$ та $p_\parallel$. Саме тому кінетичний опис потрібний там, де
анізотропія розподілу є істотною.

Із моментів можна побудувати також густину кінетичної енергії частинок,

\begin{equation}\label{eq:kinetic-energy-density}
 w_s=\frac{m_s}{2}\int v^2 f_s\,d^3v.
\end{equation}

Після розкладання $\vect v=\vect v_s+\vect c$ вона містить дві частини: енергію
впорядкованого руху та енергію хаотичного руху. Саме друга частина пов'язана
з~температурою і~тиском.



% --------------------------------------------------------
\section{Кінетичне рівняння}
% --------------------------------------------------------



Еволюцію функції розподілу визначає кінетичне рівняння. У~беззіткнювальній плазмі воно має
вигляд рівняння Власова:

\begin{equation}\label{eq:vlasov}
\frac{\partial f_s}{\partial t}
+\vect v\cdot\nabla f_s
+\frac{q_s}{m_s}
\left(
\vect E+\frac1c\vect v\times\vect B
\right)
\cdot\nabla_{\vect v}f_s=0.
\end{equation}

Перший доданок описує зміну розподілу з~часом, другий~--- його перенесення в координатному
просторі, а~третій~--- зміну швидкостей під дією сили Лоренца.

Рівняння можна записати компактніше через повну похідну вздовж траєкторії частинки:

\begin{equation}\label{eq:vlasov-characteristic}
\frac{d f_s}{dt}=0.
\end{equation}

Це не означає, що $f_s$ не змінюється в~просторі або в~часі. Йдеться про те, що вздовж
фазової траєкторії, визначеної рівняннями руху частинки, значення функції розподілу
зберігається. Фізично це є законом збереження числа частинок у~фазовому просторі.

Якщо зіткнення важливі, рівняння набуває форми

\begin{equation}\label{eq:kinetic-collision}
\frac{\partial f_s}{\partial t}
+\vect v\cdot\nabla f_s
+\frac{q_s}{m_s}
\left(
\vect E+\frac1c\vect v\times\vect B
\right)
\cdot\nabla_{\vect v}f_s
=C_s[f].
\end{equation}

Оператор $C_s[f]$ описує зміну функції розподілу внаслідок зіткнень. Для слабкозв'язаної
плазми кулонівська взаємодія далекодіюча, тому проста модель зіткнення двох твердих кульок
не є адекватною: зіткнення визначається великою кількістю слабких кулонівських відхилень.

Таким чином, рівняння Власова та кінетичне рівняння зіткнювальної плазми відрізняються
саме способом урахування взаємодії між частинками. У~першому випадку основною є
самоузгоджена електромагнітна взаємодія, у~другому додається кінетичний ефект зіткнень.



% --------------------------------------------------------
\section{Максвеллівський розподіл і~температура}
% --------------------------------------------------------



Найпростішим прикладом функції розподілу є локальний максвеллівський розподіл. Для
частинок маси $m$, концентрації $n$ та температури $T$ він має вигляд

\begin{equation}\label{eq:maxwell-distribution}
 f_M(\vect v)=
 n\left(\frac{m}{2\pi k_{\mathrm B}T}\right)^{3/2}
 \exp\left[
 -\frac{m(\vect v-\vect u)^2}{2k_{\mathrm B}T}
 \right].
\end{equation}

Тут $\vect u$ є середньою швидкістю. Розподіл має максимум при $\vect v=\vect u$,
а~його ширина визначається температурою. Чим вища $T$, тим більша частка частинок має
великі теплові швидкості.

Для максвеллівського розподілу

\begin{equation}\label{eq:thermal-speed}
 \left\langle c^2\right\rangle
 =\frac{3k_{\mathrm B}T}{m},
\end{equation}

тому характерну теплову швидкість зручно визначити як

\begin{equation}\label{eq:thermal-speed-definition}
 v_T=\sqrt{\frac{k_{\mathrm B}T}{m}}.
\end{equation}

У~плазмі електрони та іони не зобов'язані мати однакові температури:

\begin{equation*}
T_e\neq T_i.
\end{equation*}

Різні сорти частинок можуть обмінюватися енергією через зіткнення, але цей обмін не
обов'язково відбувається настільки швидко, щоб встановити спільну температуру. Тому $T_e$
та $T_i$ слід розглядати як окремі характеристики, якщо це випливає з~умов задачі.

У~слабо зіткнювальній або беззіткнювальній плазмі розподіл може істотно відхилятися від
максвеллівського. Тоді однієї температури вже недостатньо для повного опису кінетичного
стану.



% --------------------------------------------------------
\section{Від кінетики до гідродинаміки}
% --------------------------------------------------------



Рідинні рівняння можна отримати з~кінетичного рівняння, якщо замінити повну функцію
розподілу набором її моментів. Це робить особливо наочним зв'язок між кінетичним,
дворідинним та однорідинним описами.

Помножимо рівняння Власова на одиницю та проінтегруємо за швидкостями. За умови достатньо
швидкого спадання $f_s$ при $v\to\infty$ інтеграл від похідної за швидкістю перетворюється
на поверхневий інтеграл, який зникає. Отримуємо

\begin{equation}\label{eq:kinetic-continuity}
\frac{\partial n_s}{\partial t}
+\nabla\cdot(n_s\vect v_s)=0.
\end{equation}

Це рівняння неперервності для кожного сорту частинок.

Тепер помножимо кінетичне рівняння на $m_s\vect v$ та проінтегруємо. У~результаті дістаємо
рівняння руху

\begin{equation}\label{eq:kinetic-momentum}
m_sn_s
\left(
\frac{\partial\vect v_s}{\partial t}
+\vect v_s\cdot\nabla\vect v_s
\right)
=
q_sn_s
\left(
\vect E+\frac1c\vect v_s\times\vect B
\right)
-\nabla\cdot\mathsf P_s
+\vect R_s.
\end{equation}

Тут $\vect R_s$ описує передачу імпульсу між різними сортами частинок. У~беззіткнювальній
моделі цей член відсутній.

Рівняння показує важливу річ: електромагнітна сила діє на середній рух, але внутрішня
кінетика частинок входить через тензор тиску $\mathsf P_s$. Тому просте скалярне рівняння
стану $p=p(n,T)$ є вже додатковим припущенням, а~не прямим наслідком кінетичного рівняння.

Якщо перейти до наступного моменту, виникає рівняння для потоку енергії. У~ньому
з'являються моменти вищого порядку. Отже, виникає \emphx{ієрархія моментів}:

\begin{equation*}
f_s
\longrightarrow
(n_s,\vect v_s,\mathsf P_s)
\longrightarrow
\text{потік енергії}
\longrightarrow
\text{вищі моменти}
\longrightarrow\cdots
\end{equation*}

Ієрархія сама по собі не замикається: рівняння для моменту порядку $N$ містить момент
порядку $N+1$. Щоб отримати скінченну систему рідинних рівнянь, необхідне
\emph{замикання}. Прикладами є припущення про ізотропний тиск, рівняння стану,
адіабатичний зв'язок або задану теплопровідність.

Тому ланцюжок переходу між рівнями опису має вигляд

\begin{equation}\label{eq:kinetic-hierarchy}
\boxed{
f_s(\vect r,\vect v,t)
\longrightarrow
(n_s,\vect v_s,\mathsf P_s,\ldots)
\longrightarrow
\text{дворідинна модель}
\longrightarrow
\text{МГД}.
}
\end{equation}

Кожна стрілка означає не просто зміну позначень, а~втрату частини інформації та введення
певних припущень.



% --------------------------------------------------------
\section{Дворідинна картина та її зв'язок із МГД}
% --------------------------------------------------------



Для плазми з~електронів та іонів кінетичне рівняння записують окремо для кожного сорту, а~моменти
дають дві системи рідинних рівнянь зі~своїми $n_s$, $\vect v_s$, $\mathsf P_s$ (дворідинний опис). Їхнє
додавання й~віднімання, яке зводить дворідинну модель до однорідинної МГД (рівняння руху центра мас,
узагальнений закон Ома, струм $\vect j=en_e(\vect v_i-\vect v_e)$), виконано в~розд.~\ref{sec:mhd}; тут важливо
лише, що ця різниця швидкостей губиться при переході до однієї середньої швидкості, і~тому МГД не може описати
ефекти, чутливі до окремого руху електронів та іонів.



% --------------------------------------------------------
\section{Перший міст між кінетикою та рідинним описом: дисперсія Ленгмюра}
% --------------------------------------------------------



Рідинний вивід поправки Бома~--- Гросса $\omega^2\simeq\omega_{pe}^2+3k^2v_{Te}^2$
наведено в~§\ref{sec:langmuir-waves}; там показник $3$ виникає як адіабата
одновимірного стиснення. Кінетичний вивід (див.\ нижче, \eqref{eq:eps-kinetic})
показує, що ця формула справедлива лише при $\omega/k\gg v_{Te}$, тобто для довгих хвиль,
$k\lambda_{De}\ll1$. Поправка пов'язана з~тепловим рухом електронів, а~за
$k\lambda_{De}\gtrsim0{,}3$ виникає ще й~загасання Ландау, якого рідинна модель
не передбачає.

% --------------------------------------------------------
\section{Лінеаризація кінетичного рівняння та хвильовий відгук}
% --------------------------------------------------------

До цього моменту ми розглядали функцію розподілу як спосіб описати стан плазми та її
еволюцію в~часі. Тепер поставимо інше питання: \emph{що відбувається з~плазмою, якщо
її трохи відхилити від рівноваги?}

Це питання є природним продовженням теорії хвиль. У~рідинному описі ми вже бачили,
що мале відхилення густини, швидкості або магнітного поля може поширюватися в~плазмі
у~вигляді хвилі. У~кінетичному описі треба поставити те саме питання безпосередньо
для функції розподілу.

% --------------------------------------------------------
\subsection*{Що означає хвиля в~кінетичному описі?}
% --------------------------------------------------------

Нехай плазма перебуває в~простому однорідному рівноважному стані. Для кожного сорту
частинок його описує функція розподілу

\begin{equation}\label{eq:kinetic-equilibrium-distribution}
f_s=f_{0s}(\vect v).
\end{equation}

Ми спеціально розглядаємо однорідне рівноважний стан, тому $f_{0s}$ не залежить від
$\vect r$ та $t$. Фізично це означає, що в~середньому властивості плазми однакові
в~усіх точках простору та не змінюються з~часом.

Тепер уявімо, що цей стан трохи збурили. Наприклад, у~деякій області виникла невелика
надлишкова густина електронів. Вона створює електричне поле. Це поле змінює рух
електронів, а~змінений рух електронів змінює густину заряду і, відповідно, саме поле.

Отже, виникає замкнений ланцюг:

\begin{equation*}
\begin{gathered}
\text{збурення розподілу}
\;\longrightarrow\;
\text{збурення заряду}
\;\longrightarrow\;
\text{електричне поле}\\
\longrightarrow\;
\text{зміна руху частинок}
\;\longrightarrow\;
\text{нове збурення}.
\end{gathered}
\end{equation*}

Саме ця \emph{самоузгодженість} є суттю колективної хвилі в~плазмі. Поле не просто
задане ззовні: воно створюється самими частинками і~водночас впливає на їхній рух.

У~кінетичному описі тому хвиля є збуренням самої функції розподілу. Запишемо

\begin{equation}\label{eq:kinetic-wave-perturbation}
f_s(\vect r,\vect v,t)
=
f_{0s}(\vect v)
+
\delta f_s(\vect r,\vect v,t),
\qquad
|\delta f_s|\ll f_{0s}.
\end{equation}

Тут $f_{0s}$ описує незбурений стан, а~$\delta f_s$ --- малу зміну цього стану.
На відміну від густини або середньої швидкості, збурення $\delta f_s$ залежить ще й~від
швидкості $\vect v$. Тому воно відповідає не просто зміні кількості частинок, а~зміні
самої форми розподілу за швидкостями.

Це принципово важливо. Дві ділянки плазми можуть мати однакову густину, але різний
розподіл частинок за швидкостями. Кінетична теорія дозволяє розрізнити такі стани,
тоді як рідинний опис може їх ототожнювати.

% --------------------------------------------------------
\subsection*{Лінеаризація рівняння Власова}
% --------------------------------------------------------

Щоб знайти, як розвивається мале збурення, підставимо \eqref{eq:kinetic-wave-perturbation}
у~рівняння Власова \eqref{eq:vlasov}. Для простоти спочатку розглянемо електростатичні
збурення однорідної плазми без зовнішнього магнітного поля. Тоді

\begin{equation}\label{eq:kinetic-electrostatic-vlasov}
\frac{\partial f_s}{\partial t}
+
\vect v\cdot\nabla f_s
+
\frac{q_s}{m_s}\vect E\cdot\nabla_{\vect v}f_s
=0.
\end{equation}

У~рівноважному стані поле відсутнє, тому

\begin{equation*}
\vect E_0=0.
\end{equation*}

Після збурення поле також є малим:

\begin{equation*}
\vect E=\delta\vect E.
\end{equation*}

Підставимо ці вирази у~\eqref{eq:kinetic-electrostatic-vlasov}. Виникають члени двох
типів. Частина є лінійною за малими величинами, а~частина містить добуток двох
збурень, наприклад

\begin{equation*}
\delta\vect E\cdot\nabla_{\vect v}\delta f_s.
\end{equation*}

Оскільки і~$\delta\vect E$, і~$\delta f_s$ малі, їхній добуток є величиною другого
порядку мализни. У~лінійному наближенні такими членами нехтуємо. Залишається

\begin{equation}\label{eq:kinetic-linearized-vlasov}
\frac{\partial\delta f_s}{\partial t}
+
\vect v\cdot\nabla\delta f_s
+
\frac{q_s}{m_s}
\delta\vect E\cdot\nabla_{\vect v}f_{0s}
=0.
\end{equation}

Тепер маємо перше принципово нове для нас рівняння: воно безпосередньо визначає еволюцію
малого збурення функції розподілу.

У~ньому добре видно фізику процесу. Перші два доданки описують перенесення збурення
частинками. Останній доданок показує дію електричного поля: поле змінює швидкості
частинок, а~отже, перерозподіляє їх у~просторі швидкостей.

Таким чином, \eqref{eq:kinetic-linearized-vlasov} вже містить обидві необхідні складові
кінетичної картини: \emph{вільний рух частинок} та \emph{їхню реакцію на збурене поле}.

Але поки що це ще не хвильове рівняння.

% --------------------------------------------------------
\subsection*{Як із цього виникає хвиля?}
% --------------------------------------------------------

Щоб побачити хвильову поведінку, треба згадати, що збурення функції розподілу саме
створює електричне поле. З~функції $\delta f_s$ отримуємо збурення густини заряду:

\begin{equation}\label{eq:kinetic-perturbed-charge}
\delta\rho
=
\sum_s q_s
\int\delta f_s,d^3v.
\end{equation}

А~це поле визначається рівняннями електродинаміки. У~найпростішій електростатичній
задачі достатньо рівняння Пуассона:

\begin{equation}\label{eq:kinetic-perturbed-poisson}
\nabla\cdot\delta\vect E
=
4\pi\delta\rho.
\end{equation}

Отже, маємо замкнений ланцюг:

\begin{equation*}
\delta f_s
\;\longrightarrow\;
\delta\rho
\;\longrightarrow\;
\delta\vect E
\;\longrightarrow\;
\delta f_s.
\end{equation*}

Саме тому збурення не можна розглядати як просто задану зміну функції розподілу.
Воно повинно бути узгоджене з~полем, яке саме це збурення породжує.

Якщо таке самоузгоджене збурення поширюється в~просторі та змінюється в~часі, ми
називаємо його хвилею.

Це важливий момент: у~кінетичній теорії хвиля не вводиться окремим хвильовим рівнянням.
Вона виникає як \emph{самоузгоджений розв'язок лінеаризованого кінетичного рівняння
разом із рівняннями електромагнітного поля}.

% --------------------------------------------------------
\subsection*{Нормальна мода}
% --------------------------------------------------------

Тепер скористаємося спеціальною властивістю обраного рівноважного стану. Воно однорідне
та стаціонарне, тому коефіцієнти лінеаризованих рівнянь не залежать від $\vect r$ і~$t$.
Для такої системи природно розкласти довільне мале збурення на гармонічні складові.

Розглянемо одну таку складову та запишемо

\begin{equation}\label{eq:kinetic-wave-harmonic}
\delta f_s(\vect r,\vect v,t)
=\widetilde f_s(\vect v)
\exp\left[i(\vect k\cdot\vect r-\omega t)\right].
\end{equation}

Тепер ця формула вже не з'являється нізвідки. Ми не стверджуємо, що функція розподілу
\enquote{зобов'язана} мати такий вигляд. Ми розкладаємо мале збурення на нормальні
гармонічні моди, кожна з~яких характеризується хвильовим вектором $\vect k$ та частотою
$\omega$.

Для електричного поля записуємо аналогічно

\begin{equation*}
\delta\vect E(\vect r,t)
=
\widetilde{\vect E}
\exp\left[i(\vect k\cdot\vect r-\omega t)\right].
\end{equation*}

Спільний множник у~цих виразах описує просторово-часову структуру хвилі, тоді як
$\widetilde f_s(\vect v)$ визначає, як ця хвиля змінює розподіл частинок за швидкостями.

Після такої підстановки диференціювання стає простим:

\begin{equation*}
\frac{\partial}{\partial t}
\longrightarrow
-i\omega,
\qquad
\nabla
\longrightarrow
i\vect k.
\end{equation*}

Тому \eqref{eq:kinetic-linearized-vlasov} перетворюється на

\begin{equation}\label{eq:kinetic-wave-response}
-i(\omega-\vect k\cdot\vect v)\widetilde f_s
+
\frac{q_s}{m_s}
\widetilde{\vect E}\cdot
\nabla_{\vect v}f_{0s}
=0.
\end{equation}

Це рівняння вже значно відрізняється від звичайного рівняння для рідинної швидкості.
Воно містить $\vect v$ як незалежну змінну. Тому різні групи частинок із різними
швидкостями мають різний відгук на ту саму хвилю.

% --------------------------------------------------------
\subsection*{Звідки береться резонанс?}
% --------------------------------------------------------

Саме тут уперше з'являється характерна для кінетичної теорії фізика.

У~\eqref{eq:kinetic-wave-response} присутня комбінація

\begin{equation}\label{eq:kinetic-wave-resonance-denominator}
\omega-\vect k\cdot\vect v.
\end{equation}

Розглянемо для простоти рух уздовж напрямку поширення хвилі. Тоді
$\vect k\cdot\vect v=k_\parallel v_\parallel$, і~особливо важливими стають частинки,
для яких

\begin{equation}\label{eq:kinetic-wave-resonance}
v_\parallel\simeq\frac{\omega}{k_\parallel}.
\end{equation}

Права частина є фазовою швидкістю хвилі вздовж її напрямку поширення. Отже, резонансні
частинки рухаються приблизно з~тією самою швидкістю, з~якою рухається фаза хвилі.

Фізично це зрозуміло. Якщо частинка рухається значно повільніше або значно швидше за
хвилю, вона швидко проходить крізь різні фази хвилі. Вплив поля на неї при цьому
значною мірою взаємно компенсується. Але якщо частинка рухається приблизно разом
із фазою, вона може протягом тривалого часу залишатися в~області поля одного знака
і~ефективно обмінюватися з~хвилею енергією.

Таким чином, резонанс не є додатковою умовою, яку ми просто постулюємо. Він безпосередньо
випливає зі структури кінетичного відгуку \eqref{eq:kinetic-wave-response}.

Саме цієї інформації не було в~простому рідинному описі. Там ми працювали із середньою
швидкістю, густиною та тиском. Тут же бачимо, що частинки з~різними $v_\parallel$
взаємодіють із хвилею по-різному.

% --------------------------------------------------------
\subsection*{Самоузгодженість і~дисперсійне рівняння}
% --------------------------------------------------------

До цього моменту ми фактично розглядали відгук частинок на збурене електричне поле.
Але поле не є незалежною величиною: воно створюється самою плазмою. Тому треба одночасно
виконати рівняння Пуассона.

Для гармонічної моди \eqref{eq:kinetic-wave-harmonic} воно має вигляд

\begin{equation}\label{eq:kinetic-wave-poisson-harmonic}
i\vect k\cdot\widetilde{\vect E}
=
4\pi
\sum_s q_s
\int\widetilde f_s(\vect v)\,d^3v.
\end{equation}

Отже, маємо дві частини однієї задачі:

\begin{enumerate}
\item кінетичне рівняння визначає, як частинки реагують на поле;
\item рівняння Пуассона визначає, яке поле створюють ці частинки.
\end{enumerate}

Потрібно знайти такий розподіл частинок і~таке поле, які одночасно задовольняють обидві
умови. Саме ця вимога називається \emph{самоузгодженістю}.

Після підстановки кінетичного відгуку в~рівняння поля отримуємо умову існування
нетривіального розв'язку. У~загальному вигляді її записують як дисперсійне рівняння

\begin{equation}\label{eq:kinetic-wave-dispersion-general}
D(\omega,\vect k)=0.
\end{equation}

Це рівняння відповідає на два принципові питання.

По-перше, які хвилі взагалі можуть існувати в~даній плазмі? Відповідь задається
дисперсією $\omega(\vect k)$.

По-друге, що відбувається з~амплітудою хвилі під час її поширення? Якщо $\omega$ має
уявну частину, амплітуда може зростати або спадати з~часом.

Отже, дисперсійне рівняння є результатом не просто математичної підстановки. Воно є
умовою самоузгодженого існування колективного руху частинок і~поля.

% --------------------------------------------------------
\subsection*{Чому кінетичний опис дає нову фізику?}
% --------------------------------------------------------

У~рідинному описі ми замінюємо повну функцію розподілу її моментами:
густиною, середньою швидкістю, тиском тощо. Це дуже корисно, але разом із цим
втрачається інформація про те, скільки частинок має ту чи іншу швидкість.

Для багатьох повільних і~великомасштабних процесів ця втрата інформації неістотна.
Але для хвиль вона може стати принциповою. Як показує \eqref{eq:kinetic-wave-response},
відгук залежить безпосередньо від $\vect v$. Тому хвиля може вибірково взаємодіяти
з~певною частиною розподілу.

У~цьому полягає одна з~головних переваг кінетичного опису: він дозволяє досліджувати
не лише рух середнього потоку, а~й обмін енергією між хвилею та частинками з~різними
швидкостями.

Саме тут природно виникає загасання Ландау. Якщо швидкості резонансних частинок
попадають на ділянку розподілу, де кількість частинок із різними швидкостями неоднакова,
обмін енергією між хвилею та частинками не компенсується. Для звичайного
максвеллівського розподілу це приводить до загасання хвилі.

Таким чином, наступний крок буде вже не введенням нового математичного апарату заради
самого апарату. Ми використаємо отриману картину, щоб зрозуміти, \emph{чому хвиля може
втрачати енергію навіть у~повністю беззіткнювальній плазмі}.




% --------------------------------------------------------
\section{Загасання Ландау}\label{sec:landau-damping}
% --------------------------------------------------------

Одним із найважливіших результатів кінетичної теорії є можливість загасання хвилі навіть
у~беззіткнювальній плазмі. Покажемо це кількісно для ленгмюрівських хвиль.

\subsection*{Діелектрична проникність беззіткнювальної плазми}
Розглянемо одновимірну поздовжню хвилю $\propto e^{i(kx-\omega t)}$ вздовж $x$ і~позначимо
$F_{s}(v)=\int f_{0s}\,dv_y\,dv_z$ (так що $\int F_s\,dv=n_s$). З~\eqref{eq:kinetic-wave-response}
$\widetilde f_s=-i\,(q_s/m_s)\widetilde E\,F_s'/(\omega-kv)$. Підставляючи в~рівняння Пуассона
\eqref{eq:kinetic-wave-poisson-harmonic} і~позначаючи $ik\widetilde E=4\pi\widetilde\rho$,
отримуємо дисперсійне рівняння $\varepsilon(\omega,k)=0$ з~
% ~~~~~~~~~~~~~~~~~~~~~~~~~~~~~~~~~~~~~~~~~~~~~~~~~~~~~~~~
\begin{equation}\label{eq:eps-kinetic}
	\varepsilon(\omega,k)=1+\sum_s\frac{\omega_{ps}^2}{k\,n_s}\int\frac{F_s'(v)}{\omega-kv}\,dv.
\end{equation}
% ~~~~~~~~~~~~~~~~~~~~~~~~~~~~~~~~~~~~~~~~~~~~~~~~~~~~~~~~
Для холодної плазми $F_s=n_s\delta(v)$ маємо $\int F'/(\omega-kv)\,dv=-n_sk/\omega^2$ і~
$\varepsilon=1-\sum\omega_{ps}^2/\omega^2$~--- відомий результат~\eqref{eq:plasma-permittivity}.
Якщо ж $\omega\gg kv_{Te}$, то розкладання $1/(\omega-kv)=\omega^{-1}\sum(kv/\omega)^n$
і~інтегрування по частинах дають $\int F'/(\omega-kv)\,dv\approx-(n_sk/\omega^2)\bigl(1+3k^2v_{Ts}^2/\omega^2\bigr)$,
звідки $\varepsilon\approx1-(\omega_p^2/\omega^2)(1+3k^2v_{Te}^2/\omega^2)$ і~
$\varepsilon=0$ дає формулу Бома~--- Гросса~\eqref{eq:bohm-gross}.
Новим є те, що підінтегральний вираз має \emph{полюс} при $v=\omega/k$: саме тут
«живуть» резонансні частинки.

Як обійти полюс? Л.~Ландау (1946~р.) показав, що правильно розв'язувати задачу початкових
умов: тоді збурення вмикається в~момент $t=0$, і~для збіжності інтеграла треба вважати
$\operatorname{Im}\omega>0$. Результат продовжують аналітично на область
$\operatorname{Im}\omega\le0$, що зводиться до заміни $\omega\to\omega+i0$. За формулою
Сохоцького
$1/(\omega-kv+i0)=\mathrm P\,1/(\omega-kv)-i\pi\delta(\omega-kv)$ маємо
$\varepsilon=\varepsilon_r+i\varepsilon_i$ з~
% ~~~~~~~~~~~~~~~~~~~~~~~~~~~~~~~~~~~~~~~~~~~~~~~~~~~~~~~~
\begin{equation}\label{eq:eps-imag}
	\varepsilon_i(\omega,k)=-\pi\sum_s\frac{\omega_{ps}^2}{k^2\,n_s}\,F_s'\!\left(\frac{\omega}{k}\right).
\end{equation}
% ~~~~~~~~~~~~~~~~~~~~~~~~~~~~~~~~~~~~~~~~~~~~~~~~~~~~~~~~
Уявна частина проникності визначається нахилом розподілу \emph{у~точці резонансу}
$v=v_\varphi=\omega/k$~--- вона відсутня в~будь-якому рідинному описі.

\subsection*{Декремент}
Нехай хвиля слабко загасає: $\omega=\omega_r+i\gamma$, $|\gamma|\ll\omega_r$. Тоді
$\varepsilon_r(\omega_r,k)=0$ (це дає частоту, наприклад Бома~--- Гросса), а~поправка
знаходиться з~розкладу: $\gamma=-\varepsilon_i\Big/\dfrac{\partial\varepsilon_r}{\partial\omega}$.
Для ленгмюрівської хвилі $\partial\varepsilon_r/\partial\omega\approx2/\omega_{pe}$, тому
% ~~~~~~~~~~~~~~~~~~~~~~~~~~~~~~~~~~~~~~~~~~~~~~~~~~~~~~~~
\begin{equation}\label{eq:landau-general}
	\gamma=\frac{\pi\omega_{pe}^3}{2k^2n_e}\,F_e'(v_\varphi).
\end{equation}
% ~~~~~~~~~~~~~~~~~~~~~~~~~~~~~~~~~~~~~~~~~~~~~~~~~~~~~~~~
Знак $\gamma$ збігається зі~знаком $F_e'(v_\varphi)$. Для максвеллівського розподілу
$F_e'(v)=-(v/v_{Te}^2)F_e$, тобто $F_e'(v_\varphi)<0$, і~хвиля загасає:
% ~~~~~~~~~~~~~~~~~~~~~~~~~~~~~~~~~~~~~~~~~~~~~~~~~~~~~~~~
\begin{equation}\label{eq:landau-gamma}
	\tcbhighmath{\frac{\gamma}{\omega_{pe}}=-\sqrt{\frac\pi8}\,\frac{1}{(k\lambda_{De})^3}
	\exp\!\left(-\frac{1}{2k^2\lambda_{De}^2}-\frac32\right)}\qquad(k\lambda_{De}\ll1).
\end{equation}
% ~~~~~~~~~~~~~~~~~~~~~~~~~~~~~~~~~~~~~~~~~~~~~~~~~~~~~~~~
Тут $\lambda_{De}=v_{Te}/\omega_{pe}$ ($v_{Te}^2=k_{\mathrm B}T_e/m_e$), а~$-3/2$ в~експоненті
враховує поправку Бома~--- Гросса до $v_\varphi$. Експоненціальна залежність від
$1/k^2\lambda_{De}^2$ показує, що довгі хвилі ($k\lambda_{De}\ll1$) практично не
загасають, бо $v_\varphi\gg v_{Te}$ і~в~«хвості» розподілу майже немає частинок. Типові числа:
$k\lambda_{De}=0{,}3$ дає $|\gamma|/\omega_{pe}\approx0{,}02$, а~$k\lambda_{De}=0{,}5$~---
$\approx0{,}15$ (така хвиля вже живе лише кілька періодів, і~розклад за малим $\gamma$ втрачає силу).

\begin{figure}[h!]\centering
	\localinput{landau.tikz}
	\caption{Загасання Ландау: частинок, трохи повільніших за хвилю, більше, ніж трохи
	швидших, тому обмін енергією не компенсується, і~хвиля віддає енергію розподілу}
	\label{pic:landau}
\end{figure}

\subsection*{Фізична картина}
Частинки, трохи повільніші за хвилю, прискорюються нею (хвиля «наздоганяє» їх і~штовхає),
а~трохи швидші~--- гальмуються (вони передають енергію хвилі). Оскільки для спадного
розподілу ($F'<0$) повільних частинок більше, ніж швидких, сумарно енергія переходить
від хвилі до частинок (рис.~\ref{pic:landau}). Це загасання принципово відрізняється від
омічної дисипації: зіткнень немає, а~енергія перекачується з~поля в~кінетичну енергію
резонансних частинок.

\paragraph{Нестійкість «пучка на хвості».}
Якщо розподіл має локальний максимум при $v_b>v_{Te}$ (електронний пучок), то на
його висхідній ділянці $F_e'(v_\varphi)>0$, і~за~\eqref{eq:landau-general} $\gamma>0$~---
хвиля \emph{зростає} (обернене загасання Ландау, розд.~\ref{sec:instability}). Нестійкість
зникає, коли пучок розпливається або хвиля захоплює резонансні частинки.

\paragraph{Нелінійне захоплення.}
Лінійна теорія працює, поки амплітуда мала. Резонансні частинки в~потенціалі хвилі
$\varphi_0$ здійснюють коливання з~частотою відбиття $\omega_b=\sqrt{ek^2\varphi_0/m_e}$.
Коли $\omega_b\gtrsim|\gamma|$, частинки встигають повернути й~віддати енергію назад:
розподіл у~околі $v_\varphi$ вирівнюється (плато), $F'(v_\varphi)\to0$, і~загасання
припиняється. Це явище, яке експериментально підтвердили Малмберг і~Вхартон (1964~р.).

Тому загасання Ландау є важливою межею простого рідинного опису. Скалярні величини $n$,
$\vect v$ та $T$ не містять достатньої інформації про поведінку розподілу поблизу
резонансної швидкості. Потрібна сама функція $f_s(\vect r,\vect v,t)$ або принаймні
модель, яка зберігає відповідну кінетичну інформацію.



% --------------------------------------------------------
\section{Коли потрібен кінетичний опис?}
% --------------------------------------------------------



Вибір між кінетичною та рідинною моделлю визначається не тим, яка модель
\enquote{правильна}, а~масштабом явища. Якщо нас цікавлять довгохвильові процеси, для яких
розподіл частинок залишається близьким до локальної рівноваги, можна перейти до моментів
і~замкнути систему рідинними рівняннями.

Кінетичний опис стає необхідним, коли важливими є:

\begin{itemize}
\item детальний розподіл частинок за швидкостями;
\item резонансна взаємодія частинок із хвилями;
\item анізотропія розподілу або тиску;
\item нерівноважні хвости функції розподілу;
\item процеси, для яких зіткнення занадто рідкі, щоб підтримувати локальну рівновагу.
\end{itemize}

У~магнетизованій плазмі додатково виникає відмінність між рухом уздовж та поперек
магнітного поля. Тому навіть тиск може стати анізотропним,

\begin{equation*}
p_\perp\neq p_\parallel,
\end{equation*}

і~проста однорідинна МГД-модель зі скалярним тиском перестає бути достатньою.

Загальну послідовність переходів між рівнями опису, $f_s\to(n_s,\vect v_s,\mathsf P_s)\to$ дворідинна модель $\to$ МГД, наведено в~\eqref{eq:kinetic-hierarchy}.

Чим далі ми рухаємося вправо, тим менше мікроскопічної інформації зберігаємо і~тим більше
припущень про замикання повинні зробити. Натомість рідинний опис стає значно простішим для
довгохвильових макроскопічних процесів.



%% --------------------------------------------------------
\section*{Задачі}
%% --------------------------------------------------------



\begin{problem}\label{prb:KineticMoments}
Для однорідної беззіткнювальної плазми покажіть, що інтегрування рівняння Власова за всіма
швидкостями приводить до рівняння неперервності
\begin{equation*}
\frac{\partial n_s}{\partial t}
+\nabla\cdot(n_s\vect v_s)=0.
\end{equation*}
Яку властивість функції розподілу потрібно використати, щоб поверхневий інтеграл
у~просторі швидкостей зник?
\end{problem}

\begin{problem}\label{prb:MaxwellPressure}
Для максвеллівського розподілу покажіть, що
\begin{equation*}
p_s=\frac{m_s}{3}\int c^2f_s\,d^3v
=n_sk_{\mathrm B}T_s.
\end{equation*}
Поясніть, чому цей результат потребує ізотропності розподілу.
\end{problem}

\begin{problem}\label{prb:LandauResonance}
Для ленгмюрівської хвилі в~плазмі з~$T_e=10$~еВ, $n_e=10^{12}~\text{см}^{-3}$ та~довжиною хвилі
$\lambda=2\pi/k=1$~мм знайдіть $v_\varphi$, відношення $v_\varphi/v_{Te}$ і~декремент $\gamma/\omega_{pe}$
за~\eqref{eq:landau-gamma}. Чи виконується умова слабкого загасання? Поясніть, чому саме
частинки з~$v\simeq v_\varphi$ відіграють особливу роль.
\end{problem}

\begin{problem}\label{prb:BumpOnTail}
Розподіл електронів складається з~максвеллівського «тіла» ($n_0$, $v_{Te}$) та~слабкого пучка
($n_b\ll n_0$, швидкість $v_b\gg v_{Te}$, розкид $\Delta v\ll v_b$). Використовуючи
\eqref{eq:landau-general}, оцініть інкремент нестійкості $\gamma\sim\omega_{pe}(n_b/n_0)(v_b/\Delta v)^2$ та поясніть, чому
він максимальний при $v_\varphi$ трохи менше $v_b$.
\end{problem}



%% --------------------------------------------------------
\section*{Контрольні питання}
%% --------------------------------------------------------



\begin{enumerate}
\item Що таке функція розподілу $f_s(\vect r,\vect v,t)$ і~чим вона відрізняється від
  концентрації $n_s$?
\item Як через моменти функції розподілу визначають концентрацію, середню швидкість та
  тензор тиску?
\item Чому в~загальному випадку тиск є тензором, а~не скаляром?
\item Який фізичний зміст рівняння Власова?
\item Що змінюється в~кінетичному рівнянні, якщо зіткнення не можна нехтувати?
\item Чому кулонівські зіткнення не можна просто моделювати як зіткнення твердих кульок?
\item Що визначає ширину максвеллівського розподілу?
\item Чому в~плазмі можуть бути різні температури електронів та іонів?
\item Як із кінетичного рівняння отримується рівняння неперервності?
\item Чому ієрархія моментів не замикається сама?
\item Які припущення можуть використовуватися для замикання рідинної моделі?
\item У~чому фізичний зміст формули Бома~--- Гросса?
\item Чим загасання Ландау відрізняється від омічної дисипації?
\item Чому для опису загасання Ландау недостатньо знати лише $n$, $\vect v$ та $T$?
\item Від чого залежить знак декремента $\gamma$ в~формулі~\eqref{eq:landau-general}? Коли хвиля зростає?
\item Що таке нелінійне захоплення частинок хвилею і~чому воно припиняє загасання?
\item За яких обставин кінетичний опис є принципово необхідним?
\end{enumerate}
